\documentclass[12pt]{article}

\usepackage[margin=2cm]{geometry}
\usepackage{amsmath}
\usepackage{multicol}
\usepackage{nopageno}
\usepackage{SIunits}
\usepackage{graphicx}
\usepackage{enumerate}
\usepackage{enumitem}
%\usepackage{mathabx}
\usepackage{wasysym}


% for making bond drawings
\usepackage{chemfig}
% Global bond settings
\setchemfig{
  atom sep=18pt,               % fixed bond length
  double bond sep=2pt,       % spacing between the two bars
  bond style={line width=0.8pt}
}


% for centering figures on page while in an enumerated/itemized list:
\usepackage{changepage} % in the preamble
\makeatletter
\newenvironment{centerinpage}
  {%
    \par\noindent
    \hspace*{-\@totalleftmargin}% back to page margin (covers all nested list indents)
    \begin{minipage}{\textwidth}% fixed to page width (ignores widened \linewidth)
      \centering
  }
  {%
    \end{minipage}\par
  }
\makeatother

%%%%%%%%%%%%%%%%%%%%%%%%%%%%%%%%%%%%%%%%%%%%%%%%%%%%%
% For putting links in the document (both to sections and hyperlinks)
\usepackage{hyperref}
\usepackage{url}
\hypersetup{
    colorlinks=true,     % false: boxed links; true: colored links
    linkcolor=red,       % color of internal links
    citecolor=blue,      % color of links to bibliography
    filecolor=blue,      % color of file links
    urlcolor=blue        % color of external links
}
% below, use \url{http://www.nytimes.com} (they will be urlcolor)
%%%%%%%%%%%%%%%%%%%%%%%%%%%%%%%%%%%%%%%%%%%%%%%%%%%%%


%%%%%%%%%%%%%%%%%%%%%%%%%%%%%%%%%%%%%%%%%%%%%%%%%%%%%
% for the events page boxes

% boxed regions
\usepackage{tcolorbox}
\tcbuselibrary{skins,breakable}

% fonts & spacing
\usepackage{lmodern}
\usepackage{microtype}
\usepackage{calc}


% small helper macros (adjust these)
\newcommand{\boxheight}{4.2cm}   % height for each event box
\newcommand{\boxfont}{\scriptsize} % change to \tiny or \footnotesize as desired
\newcommand{\linelen}{0.95\linewidth} % timeline line length

%%%%%%%%%%%%%%%%%%%%%%%%%%%%%%%%%%%%%%%%%%%%%%%%%%%%%
\newcommand{\rmd}{\ensuremath{{\rm d}}}
%\newcommand{\vv}[1]{\ensuremath{\vec{#1}}}
\newcommand{\vv}[1]{\ensuremath{\mathbf{#1}}}
\newcommand{\ee}[1]{\ensuremath{\mathbf{e}_{#1}}}
%%%%%%%%%%%%%%%%%%%%%%%%%%%%%%%%%%%%%%%%%%%%%%%%%%%%%



\begin{document}

\begin{center}
{\Large The atmosphere and the radiative-convective equilibrium}
\end{center}
%%%%%%%%%%%%%%%%%%%%%%%%%%%%%%%%%%%%%%%%%%%%%%%%
%%%%%%%%%%%%%      overview         %%%%%%%%%%%%%%%%%%%%%
%%%%%%%%%%%%%%%%%%%%%%%%%%%%%%%%%%%%%%%%%%%%%%%%

\noindent The layers of the atmosphere, increasing from ground level (height, $h = \unit{0}{\kilo\meter}$) are
	%
	\begin{itemize}
	\item {\bf Troposphere} [0--\unit{10}{\kilo\meter}]: Region of weather, contains 80\% of total air mass.
	\item {\bf Stratosphere} [10--\unit{50}{\kilo\meter}] Ozone region where near-UV solar radiation is absorbed; temperatures increase to 0$.\degree$C
	\item {\bf Mesosphere} [50--\unit{85}{\kilo\meter}] Low density of air (and ozone) causes temperature decrease to $-100\degree$C.
	\item {\bf Thermosphere} [$>\unit{85}{\kilo\meter}$] Ionization of molecules by far-UV solar radiation causes temperatures to reach 500--2000$\degree$C but pressure and density is extremely low (the ``ionosphere'' is here).
	\end{itemize}

%
%\newpage
%%%%%%%%%%%%%%%%%%%%%%%%%%%%%%%%%%%%%%%%%%%%%%%%
%%%%%%%%%%%%%      questions         %%%%%%%%%%%%%%%%%%%%%
%%%%%%%%%%%%%%%%%%%%%%%%%%%%%%%%%%%%%%%%%%%%%%%%

\begin{enumerate}

\item Let's draw a graph of height vs.\ temperature for the lower levels of the atmosphere. Take a regular sheet of graph paper in portrait orientation. Let the vertical axis go between 0 and \unit{100}{\kilo\meter}, and the horizontal between $-60$ and 10$\degree$C.  
	%
	\begin{enumerate}
	\item Put a dot at $h=0$ and $T=15\degree$C (this will be just off the end of the lattice of the graph paper).  This represents the Earth's surface at average temperature.
	\item It can be shown with relatively simple physics (undergraduate thermodynamics) that warm rising air (convection) will cool in such a way that its temperature depends linearly on height:
		%
		\begin{equation*}
		\frac{\Delta T}{\Delta h} = - \Gamma 
		\end{equation*}
		%
	where $\Gamma$ is called the {\em adiabatic lapse rate}.  It's value for dry air (with zero humidity) can be shown to be almost exactly $\unit{10}{\kelvin/\kilo\meter}$ (or, equivalently $10 \frac{\degree\text{C}}{\text{km}}$), but the average value on Earth for {\em moist} air is about  $\unit{6.5}{\kelvin/\kilo\meter}$.  Use this value to determine the temperature at the top of the troposphere (this location is called the {\em tropopause})..  Draw a dot there and use a ruler to draw the straight line decrease of temperature from surface to tropopause.
	\item Sketch the temperature above \unit{10}{\kilo\meter} using the information given above about the atmospheric layers.
	\item The atmospheric pressure decreases exponentially with height:
		%
		\begin{equation*}
		p(h) = p_0 \, \exp \left[ -\frac{h}{h_0} \right]
		\end{equation*}
		%
	where $p_0 = \unit{1}{atm}$ and the {\em scale height}, $h_0$, is about \unit{8}{\kilo\meter}.  We can solve the above equation for height:
		%
		\begin{equation*}
		h = h_0 \ln \left[ \frac{p_0}{p} \right]
		\end{equation*}
		%
	use this to make an axis of pressure at the right of your graph paper.  Give values of pressure on this axis at the corresponding heights where $p =  \left[ 1, \frac{1}{2}, \frac{1}{4}, \frac{1}{8}, \frac{1}{16}, \frac{1}{32}, \frac{1}{64}, \frac{1}{128} \right] \times p_0$.
	\end{enumerate}

\item {\bf Radiative-Convective Equilibrium Temperature} In prior weeks we made estimates of the Earth's equilibrium temperature.  In fact this is the driving question of the first half of the course: What do we expect to be the temperature of the Earth?  Let's review our prior results and then push forward to a new (and basically final) answer using the mechanism of energy transfer by evaporation.
	%
	\begin{enumerate}
	\item Together with your group, remind yourselves of the two calculations we did previously for estimating the Earth's surface temperature: (i) Assuming no atmosphere, only the incoming solar flux ($S = \unit{1360}{\joule/\second/\meter^2}$) and the average albedo of $\alpha=0.3$; (ii) With an atmosphere that absorbs and then re-emits all of the Earths outgoing (infrared) flux.  What were the temperatures found in each case, and how do they compare to the observed average temperature of $\sim15\degree$C?
	\item The {\em lapse rate} mentioned implies an additional flux of heat from the Earth's surface to the atmosphere.  Redo the energy flow (or flux) diagram (e.g., Fig 4.5 of our textbook), but now include an additional arrow for ``convective flux,'' $F_c$, from the Earth to the atmosphere.
	\item Write down three energy balance equations ($\text{Flux}_\text{in} = \text{Flux}_\text{out}$) related to your figure from (b): (i) for the object of the whole Earth, (ii) for the atmosphere, and (iii) for the Earth's surface.
	\item It turns out we do not need to ``solve'' these equations.  We can use the {\em adiabatic lapse rate} together with your first equation (i) to find the temperature of the Earth's surface.  Assume that the height of this atmosphere ``object'' is where the atmospheric pressure drops to half its surface value.  What surface temperature do you obtain?
	\end{enumerate}

\end{enumerate}

\newpage

%%%%%%%%%%%%%%%%%%%%%%%%%%%%%%
\mbox{}
\newpage
%%%%%%%%%%%%%%%%%%%%%%%%%%%%%%

%%%%%%%%%%%%%%%%%%%%%%%%%%%%%%%%%%%%%%%%%%%%%%%%%%%%
%%%%%%%%%%               phase diagram and psychrometric chart              %%%%%%%%%%%%%
%%%%%%%%%%%%%%%%%%%%%%%%%%%%%%%%%%%%%%%%%%%%%%%%%%%%



%%%%%%%%%%%%%%%%%%%%%%%%%%%%%%%%%%%%%%%%%%%%%%%%%%%%
%%%%%%%%%%                                  References                                %%%%%%%%%%%%%%%%
%%%%%%%%%%%%%%%%%%%%%%%%%%%%%%%%%%%%%%%%%%%%%%%%%%%%
%\noindent {\bf Data Sources}: 



\end{document}
